Next, we give an injection from $P$ to the set
\[\Canon{k}(f)\times[2L(f)]\times\{\land,\lor\}\times\{0,1,x_1,\ldots,x_n,\neg x_1,\ldots,\neg x_n\}.\]
To do this we will define an operation $\DropLeaf(\phi,i)$ that takes as input a formula $\phi$ of size $s\geq 2$ and a leaf node $i\in[2s-1]$ from $\phi$ and outputs the formula $\phi'$ given as follows. We will first describe $\phi'$ informally and then give the formal description. $\phi'$ is obtained by deleting the leaf node $i$ and making the output of the node $i_p$ that $i$ fed into simply the other node that was being fed into $i_p$.
Now, we formally describe $\phi'$. Let $i_p\in[2s-1]$ be the internal node that $i$ has an edge to in $\phi$ (we know this exists because $|\phi|\geq 2$). If needed, apply a permutation to $\phi$ so $i=2s-1$ and $i_p=2s-2$. Let $u\in[2s-3]$ be the other node in $\phi$ that feeds into $v_p$. Let $\phi'$ be the formula given by the edge set
\[E'=(E\cap([2s-3]\times[2s-3]))\cup\{(u,v_{pp}):v_p\text{ feeds into }v_{pp}\text{ in }\phi\}\]
and the labelling function $\tau'=\tau|_{[2s-3]}$.
For example if $\phi=(x_1\lor x_2)\land x_3$ and 1 was then index of the $x_1$ leaf, then $\DropLeaf(\phi,1)=x_2\land x_3$.
We show this operation in some sense commutes with formula isomorphisms.
\begin{ntclaim}[Author Claim 30]
Let $\phi$ be formula of size $s$. Let $i$ be the index a leaf node in $\phi$, and let $\sigma:[2s-3]\to[2s-3]$ be a permutation. Then there exists an integer $i'$ and a permutation $\sigma':[2s-1]\to[2s-1]$ such that
\[\DropLeaf(\sigma'(\phi),i')=\sigma\circ\DropLeaf(\phi,i).\]
\end{ntclaim}
\begin{proof}
From our definition of $\DropLeaf(\cdot,\cdot)$, we can assume without loss of generality that $i=2s-1$ and that the internal node that $i$ feeds into in $\phi$ is $i_p=2s-2$.
Then the claim follows from letting $\sigma$ be equal to $\sigma'$ on $[2s-3]$ and letting $\sigma$ be the identity on $\{2s-2,2s-1\}$ and applying the various definitions.
\end{proof}
We can also define a kind of inverse operation $\AddLeaf$ function that takes the following four inputs
\begin{itemize}
\item a formula $\phi'$ on $n$-inputs of size $s$,
\item a node $i$ in the tree given by $\phi'$,
\item a gate $O\in\{\land,\lor\}$, and
\item a leaf label $\ell\in\{0,1,x_1,\ldots,x_n,\neg x_1,\ldots,\neg x_n\}$,
\end{itemize}
and outputs the formula $\phi$ of size $s+1$ given as follows. First, we give an informal description and then given a formal definition.
Intuitively, $\AddLeaf$ is adding a new $O$-gate into $\phi'$ between the $i$-node and wherever $i$ was being output to (if $i$ has an output), whose other input is a new $\ell$-labeled leaf.
We define $\phi$ formally as follows. We will use $2s+1$ to add in our new leaf and $2s$ to add in our new gate. The edge set $E_\phi$ of $\phi$ is given by taking $E_{\phi'}$ and adding in the edges $(2s+1,2s)$ and $(i,2s)$ and then, if there is a node $i_p$ that $i$ feeds into in $\phi'$, adding in an edge $(2s,i_p)$ and removing the $(i,i_p)$ edge. The node labelling $\tau_\phi$ of $\phi$ is given by
\[\tau_\phi(i)=\begin{cases}\tau_{\phi'},&\text{if }i\in[2s-1],\\O,&\text{if }i=2s,\\\ell,&\text{if }i=2s+1.\end{cases}\]
It is easy to see that $\AddLeaf$ can reverse a $\DropLeaf(\cdot,\cdot)$ operation.
\begin{ntclaim}[Author Claim 31]
Let $\phi$ be a formula of size at least two. Let $i$ be the index of a leaf node in $\phi$. Then there exists an integer $j$, a gate $O\in\{\land,\lor\}$, and a leaf label $\ell\in\{0,1,x_1,\ldots,x_n,\neg x_1,\ldots,\neg x_n\}$ such that
\[\AddLeaf(\DropLeaf(\phi,i),j,O,\ell)=\phi.\]
\end{ntclaim}
One of the main steps in our injection will be provided by the following claim.
\begin{ntclaim}[Author Claim 32]
Let $\phi\in P$. Then there is an $i$ such that $\phi'=\DropLeaf(\phi,i)$ is a size $(L(f)+k)$ formula for computing $f$.
\end{ntclaim}
\begin{proof}
Let $\phi\in P$. Then there is some internal node $j$ in $\phi$ that takes as input a leaf node indexed by $i$ satisfying $\tau_\phi(i)=b\in\{0,1\}$. We will assume $b=0$ (the proof in the $b=1$ case is similar).
Let $\phi_j$ be the the subformula computed at node $j$, and let $O_j=\tau_\phi(j)\in\{\land,\lor\}$ be the gate label of $j$. We already know that the $i$th leaf node feeds into $\phi_j$. Let $k$ be the other node feeding into $j$, and let $\phi_k$ be the subformula computed at node $k$. We split into cases depending on whether $O_j$ is an $\land$-gate or a $\lor$-gate.
First, let us suppose $O_j$ is a $\lor$-gate. Then the formula $\phi_j$ as a function is equivalent to the formula $0\lor\phi_k$ which is equivalent as a function to $\phi_k$. Hence it follows that $\phi'=\DropLeaf(\phi,i)$ is an $(L(f)+k)$-size formula (since we removed the $i$th leaf) computing $f$ (since $\phi_j$ and $\phi_k$ compute the same function).
Now, suppose that it is a $\land$-gate. Then the output of $\phi_j$ is always zero. Since $|\phi_j|\geq 2$, there exists some subformula $\phi_2$ of $\phi_j$ of size 2 (i.e. there is some node in $\phi_v$ that has two leaves as children and $\phi_2$ is the subformula computed at that node). Since $\phi_2$ has two leaves, there exists one leaf index $i'$ such that $i'\ne i$ (i.e. this leaf node is not the $i$th leaf node we were considering before).
Then we claim that $\phi'=\DropLeaf(\phi,i')$ is an $(L(f)+k)$-size formula computing $f$. It is easy to see that $|\phi'|$ is an $(L(f)+k)$-size formula since we removed the $i'$th leaf node. To see that $\phi'$ still computes $f$, note that the 0-labeled $i$th leaf node in $\phi$ still exists in $\phi'$. If the gate node $j$ was removed by the $\DropLeaf(\cdot,\cdot)$ operation, then the output wire of $\phi_j$ that computed the 0 function in $\phi$ has been replaced by the 0-leaf $i$ in $\phi'$ which still computes the 0 function, so $\phi'$ must still compute $f$. If the gate node $v$ was not removed by the $\DropLeaf(\cdot,\cdot)$ operation, then the output corresponding gate to $v$ in $\phi'$ is still computing 0 (since it is an $\land$-gate with a 0 input), so $\phi'$ still computes $f$.
\end{proof}
Now we can finally describe the injection from $P$ to the set
\[\Canon{k}(f)\times[2L(f)+2k-1]\times\{\land,\lor\}\times\{0,1,x_1,\ldots,x_n,\neg x_1,\ldots,\neg x_n\}.\]
Given an input $\phi$ in $P$, we have by Claim 32 that there exists a least $i$-value such that $\phi'=\DropLeaf(\phi,i)$ is a formula computing $f$ of size $L(f)+k$. Let $\overline{\phi'}$ be the canonical formula isomorphic to $\phi'$. Then we have that $\overline{\phi'}\in\Canon{k}(f)$ and we have that there are permutations $\sigma$ and $\sigma'$ and an integer $i'$ such that
\begin{align*}\overline{\phi'}&=\sigma(\phi')\\&=\sigma\circ\DropLeaf(\phi,i)\\&=\DropLeaf(\sigma'(\phi),i')\end{align*}
where the last equality comes from Claim 30.
Hence, by Claim 31, we have that there exists a gate index $j\in[2L(f)+2k-1]$, a gate $O\in\{\land,\lor\}$, and a leaf label $\ell\in\{0,1,x_1,\ldots,x_n,\neg x_1,\ldots,\neg x_n\}$ such that
\[\AddLeaf(\overline{\phi'},j,O,\ell,D)=\sigma'(\phi).\]
In other words, $\AddLeaf(\overline{\phi'},j,O,\ell,D)$ outputs a formula isomorphic to $\phi$.
Thus, we set the output of the injection on input $\phi$ to be $(\overline{\phi'},j,O,\ell,D)$. The fact that this is an injection from $P$ is ensured by the fact that $\AddLeaf(\overline{\phi'},j,O,\ell,D)$ is isomorphic to $\phi$ and the fact that $P$ contains only canonical formulas.
Hence, we get that
\begin{align*}|P|&\leq|\Canon{k}(f)|\cdot 2(L(f)+k)\cdot 2\cdot(2n+1)\\&\leq O(|\Canon{k}(f)|(N+k)\log N).\end{align*}
Combining this with upper bound on $\operatorname{TestPassers}$ in terms of $P$, we get that
\[|\operatorname{TestPassers}|\leq O(|\Canon{k}(f)|\cdot 2^kN\log N).\]
\end{proof}
